\section{Feedback Regimes}
\label{sec:feedback-regimes}

The distinction between proxy and target feedback becomes more useful when combined
with the \emph{kind} of feedback available.  A short feedback horizon is not equally
valuable in every regime.  What matters is both how rapidly evidence arrives and what
that evidence establishes.

We distinguish six broad regimes.  They are not mutually exclusive.  A single system
can move through several of them as an artifact progresses from simulation to
deployment.

\paragraph{Exact symbolic feedback.}
A candidate is checked against a formally specified criterion, and acceptance or
rejection follows mechanically.  Formal theorem proving is the clearest example.
Once a mathematical statement has been faithfully represented in Lean, a proof
checker provides an unusually tight relationship between the feedback used during
search and the local target of formal validity
\parencite{hubert2025alphaproof}.  Rule-complete games have a similar structure at the
level of terminal outcomes: the rules determine whether a position is legal and who
won \parencite{silver2017alphagozero}.

This regime can support extremely large search because verification is cheap relative
to generation.  Its main residual failure lies outside the verifier: the formal
statement, game objective, or specification may fail to represent what was intended.
Thus exact verification solves
\[
    \text{``Does \(x\) satisfy specification \(S\)?''}
\]
without solving
\[
    \text{``Is \(S\) the right specification?''}.
\]

\paragraph{Executable feedback.}
The candidate is run, and behavior is compared with tests, assertions, constraints, or
measured performance.  Compilation, unit tests, property tests, exploit reproduction,
and hardware benchmarks belong here.  Execution feedback enables LLMs to debug
generated programs iteratively \parencite{chen2023selfdebug}.  Algorithm optimization can
similarly combine cheap correctness checks with less frequent direct measurements of
runtime \parencite{mankowitz2023alphadev}.

Executable feedback differs from exact symbolic feedback because the tested behavior
is normally only a sample of possible behavior.  A program can pass all available
tests while remaining insecure, brittle, or incorrect outside the tested distribution.
The feedback horizon is therefore property-specific:
\[
    \tauT^{\mathrm{tested\ property}}
    \ll
    \tauT^{\mathrm{unobserved\ property}}.
\]
There is no single feedback horizon for ``the program.''

\paragraph{Repeated empirical feedback.}
The system acts on the world and obtains direct measurements sufficiently quickly to
guide the next action.  Closed-loop chemistry and autonomous materials laboratories
are increasingly close to this regime.  AlphaFlow uses in-situ measurements to guide
subsequent experiments \parencite{volk2023alphaflow}, while A-Lab integrates experiment
selection, robotic execution, synthesis, and characterization in one autonomous loop
\parencite{szymanski2023alab}.

Here the verifier is not purely computational.  Nature performs part of the
computation.  The cycle therefore has a physical floor set by manipulation, reaction,
growth, fabrication, measurement, or other dynamics of the system being studied.
Automation can remove human handoffs, but it cannot in general make
\(\tauT\rightarrow0\).

\paragraph{Noisy statistical feedback.}
Many important outcomes are not determined by one observation.  Evidence accumulates
over repeated measurements, noisy labels, heterogeneous cases, or stochastic outcomes.
Preference learning and clinical trials are examples.  A reward model compresses many
human judgments into a rapidly queried statistical evaluator; clinical surrogate
endpoints compress later patient outcomes into earlier measurements.

In this regime, a feedback event need not be wrong in isolation.  The problem is that
its evidential relationship to the target is probabilistic:
\[
    P(T\mid F) \neq 0,1.
\]
Repeated optimization can therefore exploit stable residual error in the mapping from
\(F\) to \(T\).  Reward-model overoptimization demonstrates this directly:
independently evaluated quality can fall while the learned reward continues to rise
\parencite{gao2022overoptimization,rafailov2024direct}.  Clinical surrogate endpoints show
the same structural distinction outside AI: an earlier endpoint can be useful while
remaining an imperfect predictor of survival or quality of life
\parencite{gyawali2019accelerated,naci2017accelerated,liu2024clinicalbenefit}.

\paragraph{Sparse historical feedback.}
Some targets are observed only after long delays or rare events.  Examples include
major accidents, long-horizon institutional outcomes, macroeconomic policy effects,
and sufficiently rare safety failures.  In such settings it may be impossible to
generate repeated target-level observations on demand.

The main difficulty is therefore not merely noisy evaluation but low target-observation
rate.  If catastrophic failures occur once per \(10^6\) system-hours, a million
hours of safe operation is still only weak evidence about failure probabilities much
below that scale.  Mature safety fields respond by constructing intermediate evidence:
near misses, precursor events, stress tests, component reliability estimates,
simulation, and probabilistic risk models.  These do not eliminate the long target
horizon.  They attempt to infer the target from more frequent observations.

\paragraph{Endogenous or reflexive feedback.}
Finally, the system being optimized may influence the process that evaluates it.
Recommendation algorithms affect the users whose clicks constitute future training
data.  Economic policy changes the economy from which subsequent measurements are
taken.  A strategically capable AI may alter, anticipate, select for, or otherwise
condition on its evaluators.

The distinction is structural.  In the earlier regimes we can approximately write
\[
    \theta \longrightarrow \text{outcome}
    \longrightarrow F,
\]
with \(F\) treated as an external observation.  In a reflexive regime,
\[
    \theta
    \longrightarrow
    \{\text{outcome},\,\text{measurement process},\,\text{future data}\}
    \longrightarrow F.
\]
The evaluator is now part of the environment being optimized.

This is more demanding than ordinary proxy error.  Even an initially high-fidelity
evaluator can become less informative under optimization because the distribution or
the measurement process changes.  Real-time engagement optimization provides a mild
example: clicks can be measured almost instantly, yet optimizing them need not optimize
the user's ideal preferences or welfare \parencite{khambatta2023ideal}.  More adversarial
forms are a central concern in AI alignment, but are not yet supported by a comparably
clean empirical literature.

The feedback-horizon quantity \(K\) is most informative when read together with this
classification.  Large \(K\) under exact symbolic feedback can be benign.  Moderate
\(K\) under a systematically exploitable statistical evaluator can already be
problematic.  Under sparse or endogenous feedback, even defining an independent
target check may be the hard part.
